\documentclass[12pt]{article}
\usepackage[utf8]{inputenc}
\usepackage[a4paper, margin=2cm]{geometry}
\usepackage{amsthm}
\usepackage{amsmath, amssymb}
\usepackage{circuitikz}
\usepackage{graphicx}
\usepackage{karnaugh-map}
\usepackage{newtxtext, newtxmath}
\usepackage{tcolorbox}
\usepackage{tikz}
\usepackage{titlesec}
\usepackage{xcolor}
\usepackage{xeCJK}
\usepackage{listings}
\definecolor{lightblue}{RGB}{135, 206, 250}
\tcbuselibrary{skins,theorems}

\titleformat{\section}[block]{\Large\bfseries\centering}{}{0pt}{}
\titleformat{\subsection}[block]{\bfseries}{\thesubsection}{2pt}{}

\newtcbtheorem[no counter]{proposition}{\large\hspace{-8pt}Proposition}{
  enhanced,
  colback=white,
  colframe=lightblue,
  left=12pt,
  right=12pt,
  top=24pt,
  bottom=24pt,
  fonttitle=\bfseries
}{prop}
\newtcbtheorem[]{question}{\large\hspace{-8pt}Question :}{
  enhanced,
  colback=white,
  colframe=black,
  left=12pt,
  right=12pt,
  top=24pt,
  bottom=24pt,
  fonttitle=\bfseries
}{qst}
\newenvironment{Answer}{
  \vspace{12pt}
  {\large\textbf{\textit{\hspace{-12pt}Answer:}}}\par
  \vspace{8pt}
} {\vspace{24pt}}
\newenvironment{Proof}{
  \vspace{12pt}
  {\large\textbf{\textit{\hspace{-12pt}Proof:}}}\par
  \vspace{8pt}
} {\vspace{24pt}}

\lstset{
    language=Python,
    basicstyle=\ttfamily\footnotesize,
    keywordstyle=\color{blue}\bfseries,
    commentstyle=\color{green},
    stringstyle=\color{red},
    showstringspaces=false,
    numberstyle=\tiny\color{gray},
    stepnumber=1,
    numbersep=10pt,
    frame=single,
    breaklines=true,
    captionpos=b
}

\title{\textbf{Alogithm Design Homework 1}}
\date{May 30, 2024}

\begin{document}
\maketitle

\begin{question}{}{}
有4门选修课，学时分别为\{2, 1, 3, 2\}，学分分别为\{1.2, 1.0, 2.0, 1.5\}，
假设你有5个学时的可用时间，你会选择哪些课？可得多少学分？
给出选择策略和相应的计算过程。
\end{question}
\begin{Answer}
这是一个典型的0/1背包问题。
在这个问题中，我们需要在给定的学时内选择最大化学分的课程组合。
我们将通过动态规划方法来解决这个问题。

我们定义一个二维数组\texttt{dp[i][j]}，
其中 \(i\) 表示考虑前 \(i\) 门课，\(j\) 表示可用的学时数。
\texttt{dp[i][j]} 的值表示在这些条件下可以获得的最大学分。

动态规划的状态转移方程为：
$$dp[i][j]=\max(dp[i-1][j],\;dp[i-1][j-times[i]] + credits[i])\;\;j\geq times[i]$$

具体的 Python 代码如下：
\begin{lstlisting}
def max_credits(times, credits, max_time):
    n = len(times)
    dp = [[0] * (max_time + 1) for _ in range(n + 1)]

    for i in range(1, n + 1):
        for j in range(max_time + 1):
            if j >= times[i - 1]:
                dp[i][j] = max(dp[i - 1][j], dp[i - 1][j - times[i - 1]] + credits[i - 1])
            else:
                dp[i][j] = dp[i - 1][j]

    return dp[n][max_time]
\end{lstlisting}

通过运行上述代码，我们找到最大可得学分为 \(3.7\) 学分，选择的课程包括学时为 \(2, 1, 2\) 的课程，对应的学分为 \(1.2, 1.0, 1.5\)。
\end{Answer}

\begin{question}{}{}
有5门选修课，学时分别为\{5, 15, 22, 27, 30\}，
学分分别为\{1.2, 3.0, 4.4, 4.6, 4.7\}，
假设你有50个学时的可用时间，你会选择哪些课?可得多少学分？
给出选择策略和相应的计算过程。
\end{question}
\begin{Answer}
这是另一个典型的0/1背包问题，但与第一个问题相比，选项和总容量都有所增加。
我们需要在给定的50个学时内选择最大化学分的课程组合。
我们依旧采用动态规划方法解决这个问题。

我们定义一个二维数组 \texttt{dp[i][j]}，
其中 \(i\) 表示考虑前 \(i\) 门课，\(j\) 表示可用的学时数。
\texttt{dp[i][j]} 的值表示在这些条件下可以获得的最大学分。

动态规划的状态转移方程为：
$$dp[i][j]=\max(dp[i-1][j],\;dp[i-1][j-times[i]] + credits[i]\;\;j\geq times[i])$$

具体的 Python 代码如下：
\begin{lstlisting}
def max_credits_new(times, credits, max_time):
    n = len(times)
    dp = [[0] * (max_time + 1) for _ in range(n + 1)]

    for i in range(1, n + 1):
        for j in range(max_time + 1):
            if j >= times[i - 1]:
                dp[i][j] = max(dp[i - 1][j], dp[i - 1][j - times[i - 1]] + credits[i - 1])
            else:
                dp[i][j] = dp[i - 1][j]

    selected_courses = []
    j = max_time
    for i in range(n, 0, -1):
        if dp[i][j] != dp[i-1][j]:
            selected_courses.append(i)
            j -= times[i-1]

    selected_courses.reverse()
    return dp[n][max_time], selected_courses
\end{lstlisting}

运行上述代码后，我们发现可以获得的最大学分是9.0，
选择的课程为第3和4门课，
这些课程的学时分别为22, 27，对应的学分为4.4, 4.6。
\end{Answer}

\begin{question}{}{}
上述两个问题实例的选择策略是否相同？
给出相同或不同的理由。
\end{question}
\begin{Answer}
尽管两个问题都采用了动态规划来解决0/1背包问题，解题策略在本质上是相同的。
每个问题都定义了一个二维数组 \texttt{dp[i][j]}，
其中 \(i\) 表示考虑到的课程数量，而 \(j\) 表示可用的总学时。
动态规划的核心是使用状态转移方程来递归地计算可能的最大学分。

两个问题的解题策略相同，
都是通过比较包含当前课程与不包含当前课程的学分总和来更新 \texttt{dp} 数组。
具体而言：
$$dp[i][j]=\max(dp[i-1][j],\;dp[i-1][j-times[i]] + credits[i])\;\;j\geq times[i]$$

不同之处在于每个问题的具体参数（课程的学时和学分）以及总的学时限制不同，
这导致了不同的最优解和课程选择。
因此，虽然解题方法相同，
但每个问题的最优课程组合和最终得分可能会有显著差异。
例如，在第一个问题中，总的学时限制为5个学时，而在第二个问题中，总的学时限制为50个学时。

两个问题使用了相同的动态规划策略来解决不同的实例参数，
这展示了动态规划在处理具有不同输入但相同问题结构的情况下的适应性和强大性。
\end{Answer}

\begin{flushright}\begin{tabular}{rl}
  姓名: & 李俊洁\\
  班级: & 计算机2204 \\
  学号: & 20225443 \\
  课程: & 算法设计与分析 \\
  日期: & May 30, 2024 \\
  制作: & \LaTeX \\
  & 感谢您的批阅！
\end{tabular}\end{flushright}
\end{document}
